\chapter{The PEG Parser, Dict Union, and Pattern Matching (Python 3.9--3.10)}

The 3.9--3.10 cycle delivered the most visible syntax addition since Python 3.0: \textbf{structural
pattern matching} (\texttt{match}/\texttt{case}, PEP 634). It became possible only because 3.9
replaced the 30-year-old LL(1) parser with a \textbf{PEG parser} (PEP 617). The same window added
the daily-useful \textbf{dict union operators} \texttt{|} and \texttt{|=} (PEP 584). This chapter
treats pattern matching as the deep topic---it is \emph{destructuring plus matching}, not a C
\texttt{switch}---and is precise about the one feature that bites everyone: a bare name in a pattern
\textbf{captures}, it does not compare.

\section*{Section Index}
\begin{itemize}
  \item \textbf{15.1} The PEG parser and the syntax it unlocked (PEP 617)
  \item \textbf{15.2} Dictionary union operators \texttt{|} and \texttt{|=} (PEP 584)
  \item \textbf{15.3} Structural pattern matching (PEP 634)
  \item \textbf{15.4} The capture-vs-value gotcha, performance, and anti-patterns
  \item \textbf{15.5} Summary and cross-references
\end{itemize}

\section{The PEG parser and the syntax it unlocked (PEP 617)}

Chapter 1 covered the LL(1)$\to$PEG transition in depth; here is what it \emph{bought}. The
\textbf{PEG} parser (PEP 617, 3.9) brings \textbf{ordered choice} (first match wins) and
\textbf{unlimited lookahead} with \textbf{packrat memoization} (linear-time). Freed from one-token
lookahead, the grammar could express constructs the old parser could not---most visibly the
\textbf{parenthesized context managers} of 3.10:

\begin{lstlisting}[language=Python, caption={Parenthesized, multi-line context managers --- only parseable with PEG lookahead.}]
with (
    open("a.txt") as a,
    open("b.txt") as b,
):
    ...
\end{lstlisting}

The PEG rewrite also improved error messages and made the \texttt{match} grammar feasible. Parser
internals and the formal grammar are Vol X.

\section{Dictionary union operators \texttt{|} and \texttt{|=} (PEP 584)}

\textbf{Why this exists.} Merging dicts meant either \texttt{\{**d1, **d2\}} (cryptic) or
\texttt{m = d1.copy(); m.update(d2)} (two statements). \textbf{PEP 584} gave \texttt{dict} the binary
\texttt{|} (merge into a new dict) and \texttt{|=} (in-place update), via the
\texttt{nb\_or}/\texttt{nb\_inplace\_or} slots. Right-hand keys win.

\begin{lstlisting}[language=Python, caption={| merges (right wins), |= updates in place; non-mappings raise.}]
d1, d2 = {"a": 1, "b": 2}, {"b": 99, "c": 3}
print("d1 | d2 (right wins):", d1 | d2)

merged = d1.copy()
merged |= d2
print("|= update:", merged)

try:
    d1 | [1, 2]                  # right operand isn't a mapping
except TypeError as e:
    print("non-mapping:", str(e).split(":")[0])
\end{lstlisting}

Verified output (CPython 3.13.5):

\begin{lstlisting}[caption={Output}]
d1 | d2 (right wins): {'a': 1, 'b': 99, 'c': 3}
|= update: {'a': 1, 'b': 99, 'c': 3}
non-mapping: unsupported operand type(s) for |
\end{lstlisting}

\texttt{|} always returns a plain \texttt{dict} (even from subclasses); \texttt{|=} accepts any
mapping or key/value iterable on the right, while \texttt{|} requires a mapping.

\section{Structural pattern matching (PEP 634)}

\textbf{Why this exists --- and what it is not.} \texttt{match}/\texttt{case} is \textbf{not} a C
\texttt{switch}. It is \textbf{structural matching}: each \texttt{case} is a \emph{pattern} that
tests the \emph{shape} of the subject and \emph{binds} names from its components---closer to Rust's
\texttt{match}. It shines branching on the \emph{structure} of data (ASTs, JSON-ish messages,
command tokens).

\begin{lstlisting}[language=Python, caption={The pattern taxonomy --- literal, sequence, mapping, class, OR, guard, capture, wildcard.}]
def classify(x):
    match x:
        case 0:                              # literal (value) pattern
            return "zero"
        case [a]:                            # sequence pattern, length 1
            return f"one-list:{a}"
        case [a, b, *rest]:                  # sequence with star-capture
            return f"seq:{a},{b},rest={rest}"
        case {"type": t, "value": v}:        # mapping pattern (subset of keys)
            return f"map:{t}={v}"
        case str() as s if len(s) > 3:       # class pattern + 'as' binding + guard
            return f"long-str:{s}"
        case int() | float():                # OR pattern of class patterns
            return f"number:{x}"
        case _:                              # wildcard
            return "unknown"

for v in [0, [7], [1, 2, 3, 4], {"type": "pt", "value": 9}, "hello", 3.5, (1, 2)]:
    print(f"  {v!r:>22} -> {classify(v)}")
\end{lstlisting}

Verified output (CPython 3.13.5):

\begin{lstlisting}[caption={Output}]
                       0 -> zero
                     [7] -> one-list:7
            [1, 2, 3, 4] -> seq:1,2,rest=[3, 4]
  {'type': 'pt', 'value': 9} -> map:pt=9
                 'hello' -> long-str:hello
                     3.5 -> number:3.5
                  (1, 2) -> seq:1,2,rest=[]
\end{lstlisting}

\texttt{(1, 2)} matched the \textbf{sequence} pattern: sequence patterns match any
\texttt{collections.abc.Sequence} \emph{except} \texttt{str}/\texttt{bytes}/\texttt{bytearray}.
Mapping patterns match a \emph{subset} of keys.

\begin{lstlisting}[language=Python, caption={Class patterns destructure by \_\_match\_args\_\_ (positional) or by keyword.}]
class Point:
    __match_args__ = ("x", "y")
    def __init__(self, x, y):
        self.x, self.y = x, y

def locate(p):
    match p:
        case Point(0, 0):  return "origin"
        case Point(x, 0):  return f"x-axis@{x}"
        case Point(0, y):  return f"y-axis@{y}"
        case Point(x, y):  return f"point@{x},{y}"

print("  Point(0,0):", locate(Point(0, 0)))
print("  Point(5,0):", locate(Point(5, 0)))
print("  Point(2,3):", locate(Point(2, 3)))
\end{lstlisting}

Verified output (CPython 3.13.5):

\begin{lstlisting}[caption={Output}]
  Point(0,0): origin
  Point(5,0): x-axis@5
  Point(2,3): point@2,3
\end{lstlisting}

\textbf{What the compiler does.} A \texttt{match} compiles to a decision tree with dedicated opcodes
---\texttt{MATCH\_SEQUENCE}, \texttt{MATCH\_MAPPING}, \texttt{MATCH\_KEYS}, \texttt{MATCH\_CLASS},
plus \texttt{GET\_LEN}/\texttt{UNPACK\_SEQUENCE}---testing shape once and failing fast before binding.

\begin{lstlisting}[language=Python, caption={Real 3.13 bytecode for a sequence pattern.}]
import dis
def seqmatch(data):
    match data:
        case [1, y]:
            return y
        case _:
            return None
dis.dis(seqmatch)
\end{lstlisting}

Verified output (CPython 3.13.5, excerpt):

\begin{lstlisting}[caption={Output}]
 56   LOAD_FAST         0 (data)
      MATCH_SEQUENCE
      POP_JUMP_IF_FALSE  ... (to L1)
      GET_LEN
      LOAD_CONST        1 (2)
      COMPARE_OP        72 (==)        # length must be 2
      POP_JUMP_IF_FALSE  ... (to L1)
      UNPACK_SEQUENCE   2
      LOAD_CONST        2 (1)
      COMPARE_OP        88 (bool(==))  # first element must equal 1
      POP_JUMP_IF_FALSE  ... (to L1)
      STORE_FAST        1 (y)          # only now bind y
\end{lstlisting}

Bindings happen only on the success path---a failed sub-pattern cleans the stack, so a half-matched
pattern never leaks a partial binding.

\section{The capture-vs-value gotcha, performance, and anti-patterns}

\textbf{The one gotcha that bites everyone: a bare name captures.} An unqualified name is a
\textbf{capture pattern}---it matches \emph{anything} and binds the name; it is \textbf{not} a
comparison. To compare against a named constant, use a \textbf{dotted (value) pattern}.

\begin{lstlisting}[language=Python, caption={Bare name = capture (always matches); dotted name = value (compares).}]
import http

def capture(code):
    match code:
        case x:                       # bare name: captures, always matches
            return f"captured x={x}"
print("bare-name captures:", capture(404))

def check(code):
    match code:
        case http.HTTPStatus.OK:      # dotted: VALUE pattern, compares equality
            return "ok"
        case _:
            return "other"
print("dotted value 200:", check(200))
print("dotted value 404:", check(404))
\end{lstlisting}

Verified output (CPython 3.13.5):

\begin{lstlisting}[caption={Output}]
bare-name captures: captured x=404
dotted value 200: ok
dotted value 404: other
\end{lstlisting}

\begin{lstlisting}[language=Python, caption={3.13 rejects a capture that shadows later cases.}]
try:
    compile("match c:\n case NAME:\n  pass\n case 1:\n  pass\n", "<s>", "exec")
except SyntaxError as e:
    print("SyntaxError:", str(e).split(" (")[0])
\end{lstlisting}

Verified output (CPython 3.13.5):

\begin{lstlisting}[caption={Output}]
SyntaxError: name capture 'NAME' makes remaining patterns unreachable
\end{lstlisting}

\textbf{Other notes.} Don't use \texttt{match} for plain scalar equality (an \texttt{if/elif} or dict
dispatch is clearer); remember the str/bytes exclusion in sequence patterns; mapping patterns are
subset matches (use \texttt{**rest} or guards for exactness); \texttt{\_\_match\_args\_\_} is your
class's public destructuring contract (dataclasses set it automatically).

\section{Summary and cross-references}

\begin{itemize}
  \item The \textbf{PEG parser} (PEP 617, 3.9) replaced LL(1) and unlocked richer syntax.
  \item \textbf{\texttt{|}/\texttt{|=}} (PEP 584) merge/update dicts; \texttt{|} returns a plain
    \texttt{dict}, right keys win.
  \item \textbf{Structural pattern matching} (PEP 634) is \textbf{destructuring + matching},
    compiled to a decision tree; patterns include literal, capture, wildcard, sequence (excludes
    str/bytes), mapping (subset), class (\texttt{\_\_match\_args\_\_}), \texttt{|}, \texttt{as}, and
    guards.
  \item \textbf{A bare name captures, not compares}---use a dotted value pattern for constants; 3.13
    errors on captures that make later cases unreachable.
\end{itemize}

\textbf{Cross-references.} LL(1)$\to$PEG in depth $\rightarrow$ Chapter 1. \texttt{\{**a, **b\}}
$\rightarrow$ Chapter 11. dict internals $\rightarrow$ Vol XII. \texttt{\_\_match\_args\_\_} on
dataclasses $\rightarrow$ Chapter 13. Enums in \texttt{match} $\rightarrow$ Chapter 10. The formal
grammar $\rightarrow$ Vol X. Builtin generics (PEP 585) and \texttt{X | Y} unions (PEP 604)
$\rightarrow$ Chapter 16.
